\begin{abstract}
检索增强生成（Retrieval-Augmented Generation，RAG）把外部文档引入语言模型的生成过程，为知识更新、长尾事实问答和答案溯源提供了便利。早期系统通常对所有输入执行相同的“一次检索—生成”流程，遇到模型已经熟知的问题会浪费检索开销，面对多跳问题又难以一次找齐证据，检索结果中的噪声还可能干扰原有判断。围绕这些困难，本文从过程控制的角度梳理自适应 RAG，重点讨论四个问题：何时检索、如何形成查询、怎样评价和修正证据，以及何时停止。文章比较 FLARE、Self-RAG、Adaptive-RAG、DRAGIN、CRAG、IRCoT 等方法的控制信号与适用任务，并结合 RGB、RAGAS、ARES、RAGChecker 和 CRAG Benchmark 讨论评价方式。在方法比较之外，本文将检索视为一种有成本的知识获取行动：问题复杂度和模型置信度只能提供间接线索，证据覆盖情况及新增检索带来的收益更值得用于停止决策。基于这一认识，文章给出显式记录证据状态的设想，并提出兼顾无答案问题、噪声环境和系统开销的比较方案。

\noindent\textbf{关键词：}检索增强生成；大语言模型；自适应检索；查询改写；迭代推理；证据评价
\end{abstract}
